\renewcommand{\thesection}{S\arabic{section}}
\setcounter{section}{0}

\section{Construction of the illustrative example in Figure 1}
\label{sec:Appendix_S1}

Figure~\ref{fig:intro} provides a data-informed illustration of how
delay filtering suppresses rapid differences between candidate upstream epidemic trajectories. The figure is intended to illustrate the observation and separability framework rather than to provide a substantive reconstruction of the U.S. epidemic.

\subsection{Observed downstream incidence}

Daily U.S. confirmed-case incidence was obtained from the Johns Hopkins University Center for Systems Science and Engineering COVID-19 data stream \cite{dong2020interactive}, accessed through the Delphi Epidata API \cite{farrow2015delphi}. We used the national-level \texttt{jhu-csse/confirmed\_incidence\_num} signal. The complete series from April 1, 2020, through June 30, 2021, was used to estimate the smoothed downstream mean, construct the candidate upstream trajectories, and perform convolution, thereby preserving the pre-window history required for observations near the beginning of the displayed interval. The RMSE constraints, spectral quantities, and observation-variance reference reported in Fig.~\ref{fig:intro} were evaluated over the displayed period from September 1, 2020, through March 31, 2021.

\subsection{Smoothed downstream mean and observation noise}

To separate broad epidemic variation from day-to-day observational
fluctuation, we estimated a smooth downstream mean \(\widehat{\mu}(t)\) using second-order trend filtering. The effective degrees of freedom were selected over a prespecified grid. Conditional on each candidate smooth mean, Gaussian, Poisson, and negative-binomial observation models were fitted, and the combination of trend-filter
complexity and observation model was selected by the Akaike information criterion (AIC). The selected observation model and its fitted variance parameters were used to construct the one-sided observation-variance reference shown in Fig.~\ref{fig:intro}F, as defined below.

\subsection{Delay distribution}

The delay kernel was based on the published COVID-19
symptom-onset-to-case-report distribution reported by Zhang et al.~\cite{zhang2020evolving}. A lognormal representation was converted to a discrete daily probability mass function by integrating the fitted distribution over consecutive one-day intervals. The kernel was truncated at 60 days and renormalized to sum to one.

\subsection{Construction of the reference upstream trajectory}

The reference upstream trajectory \(U_0(t)\) was obtained by solving a nonnegative regularized inverse problem. Let \(K_g\) denote the linear convolution operator induced by the delay kernel \(g\). We estimated
\(U_0\) by minimizing
\[
\left\|K_g U-\widehat{\mu}\right\|_2^2
+
\lambda\left\|\Delta^2 U\right\|_1
\]
subject to \(U(t)\geq 0\), where \(\Delta^2\) is the second-difference operator. The regularization term stabilizes the inverse problem and produces a smooth reference trajectory whose delayed projection follows the estimated downstream mean.

\subsection{Construction of the alternative upstream trajectory}

The alternative trajectory \(U_1(t)\) was constructed to differ visibly
from \(U_0(t)\) at short temporal scales while producing a similar
downstream projection. Candidate perturbations were formed from mixtures of sinusoidal components concentrated near frequencies corresponding to 7- and 3-day periods. For each candidate perturbation \(h(t)\), we considered
\[
U_1(t)=\max\{U_0(t)+a h(t),0\}.
\]
The amplitude \(a\) was increased to the largest value satisfying both
\[
\operatorname{RMSE}(K_gU_1,K_gU_0)
\leq
\gamma\,\operatorname{RMSE}(K_gU_0,y)
\]
and
\[
\operatorname{RMSE}(K_gU_1,y)
\leq
\eta\,\operatorname{RMSE}(K_gU_0,y).
\]
For the displayed example, \(\gamma=1\) and \(\eta=1.05\). Among feasible candidates, we selected the perturbation producing the largest visible upstream separation while retaining substantial spectral energy near the two target frequencies. The RMSE constraints were evaluated over the displayed period from September 1, 2020, through March 31, 2021.

\subsection{Spectral contrast and observation-noise reference}

The spectral quantities in Fig.~\ref{fig:intro}D--F were computed over the displayed period from September 1, 2020, through March 31, 2021, using the same one-sided discrete Fourier transform and frequency grid. Panel D shows the upstream spectral contrast
\[
P_{\Delta U}(f_m)
=
\frac{q_m}{T_w}
\left|
\Delta\widetilde U_m
\right|^2,
\]
where \(\Delta\widetilde U_m\) is the discrete Fourier transform of \(U_1-U_0\) over the displayed window, \(T_w\) is the window length, and \(q_m\) is the standard one-sided multiplicity weight. Panel E shows the squared delay attenuation
\[
\left|\widetilde G_m\right|^2,
\]
and panel F shows the corresponding filtered spectral contrast
\[
p^2
\left|\widetilde G_m\right|^2
P_{\Delta U}(f_m).
\]

The horizontal line in panel F is the conservative
observation-variance normalization placed on the one-sided scale used for interior positive frequencies. For Gaussian observations, the reference level is \(2\sigma_\varepsilon^2\). For Poisson and negative-binomial observations, it is respectively
\[
2m_{\min}
\quad\text{and}\quad
2\left(
m_{\min}+\frac{m_{\min}^2}{\phi}
\right),
\]
where
\[
m_{\min}
=
\min_{i\in\{0,1\}}
\min_{t\in W}
(K_gU_i)(t)
\]
is the minimum expected downstream incidence across the two candidate trajectories over the displayed period \(W\), and \(\phi\) is the negative-binomial dispersion parameter. The factor of two accounts for the pairing of positive- and negative-frequency coefficients at interior frequencies. Panel F illustrates the frequency-wise signal term entering the separability functional, whereas the horizontal line represents its conservative variance normalization rather than a likelihood-ratio threshold.